\chapter{Outer tests and state entry when the native $r$ is infinitesimal}
\label{component:p364e14238053}
\noindent\textbf{Component }\path{p364e14238053}\par
{\small\noindent Source: \path{R13/baseline/baseline/baseline/baseline/advances/c1/centered_a_outer.tex}\par}
{\small\noindent SHA256: \path{364e1423805318f1c173683bd3eb63c27eba6442d586b66998103a7db65671ec}\par}
{\small\noindent Inventory use: CAO1--CAO8 and all state-entry paragraphs | CAO1-CAO8 and full outer/state lemma | C1\_\allowbreak{}BARE\_\allowbreak{}OTHER\_\allowbreak{}TERMINALS; precise subordinate uses retained below\par}
\medskip
\section[Outer tests and state entry when the native r is infinitesimal]{Outer tests and state entry when the native $r$ is infinitesimal}
\label{p364e14238053:sec:centered-a-outer}

\begin{lemma}
Use the unchanged physical native C1 source, normalized by $c=2$ and
$\lambda=n=0$, with all original event signs and positive openings
$a,t,M,J$. Suppose $k,\zeta,r=1-a\eta$ are positive infinitesimals and
$\eta,e=k\sqrt{\eta/\zeta}$ are positive bounded units.
Then the outer-test and state-entry arguments reduce the all-tests
problem to precisely
\[
 u,v,w>0\text{ infinitesimal},\qquad Z=1,
 \qquad w=(W+uY)/(1+u^2).
 \tag{CAO1}
\]
All other independent projective tests and states admit a full-parent
tensor SOS. No comparison of $k^2$ with $r$ is made, and arbitrarily
small positive openings and literal $\beta=0$ are included.
\end{lemma}

\paragraph{Source ratios with the common factor cancelled.}
Here $r$ always denotes the native source parameter, not the event root
$uW-Y$. Put
\[
 K=1+k^2,\quad H_c=r+\eta,\quad V_c=k^2\eta\zeta,\quad
 S_3=1-k^2\eta+tV_c,
\]
The two required short native polynomials are
\begin{align*}
P_4={}&1+k^2+\eta k^2(a-k^2-1)-\zeta k^2t\\
 &+V_c[-\zeta k^2t^2+2k^2t+a(1-t)+t],\\
P_d={}&(1+k^2)^2-a\eta k^4
 +\zeta k^2[(k^2+1)-(k^2+2)t]\\
 &+V_c[-\zeta k^2t^2+k^2a(t-1)+k^2t+k^2+1].
\end{align*}
They are positive bounded units with residue one. The exact source
identities, after cancelling the common factor $r$ in the coquotients,
are
\begin{align*}
 d&=\frac{4a\eta P_d}{KH_cS_3},&
 Ar&=\frac{2a\eta^2P_4}{KH_cS_3},\\
 \frac md&=\frac{e^2K(1-t\zeta)(1+V_c)H_c}{2M\eta P_d},&
 \frac cd&=\frac{KH_cS_3}{2a\eta P_d},\\
 \frac Ad&=\frac{\eta P_4}{2rP_d},&
 \frac Dd&=\frac{KH_c(1+V_c)-J\eta P_d}{2J\eta P_d}.
 \tag{CAO2}
\end{align*}
The positive openings obey
$M+t=1+k^2r$ and $J+a=1+k^2-k^2\zeta t$; hence $a,t,M,J$ are
bounded. Since $a\eta=1-r$, one has
\[
 \bar\eta\ge1,\qquad \bar d=4/\bar\eta,\qquad Ar\sim2,
 \qquad A\text{ positive unbounded}.
 \tag{CAO3}
\]
The numerator of $D/d$ in (CAO2) has residue
$\bar H_c-\bar J\bar\eta=1$. Therefore $m/d,c/d,D/d,A/d$ all have
positive real lower bounds. This conclusion uses no inverse bound on
an opening. Moreover,
\[
 \begin{cases}
 \displaystyle \bar D=\frac{2}{\bar\eta(\bar\eta-1)},
       &\bar\eta>1,\\[4pt]
 D\text{ positive unbounded},&\bar\eta=1.
 \end{cases}
 \tag{CAO4}
\]

\paragraph{The previously proved outer routing still applies.}
Clip $m,c,A,D$ at positive real multiples of $d$, retaining all their
nonnegative excess original event squares. The resulting coefficients
divided by $d$ are bounded. The outer-test argument in
\path{bounded_e_all_tests.tex}, through (BE2)--(BE4), uses only
these lower bounds and the independent physical-source admission at
noncommuting residue tests. Its exact homogeneous identities and
four-root ideal memberships are unchanged. It therefore includes all
projective infinity charts, the $(0,\infty)$, $(\infty,0)$ and
$(\infty,\infty)$ corners, and the unit-$W$, infinitesimal-$Z$ state
chart. The original-event SOS on literal zero axes is unchanged as
well. This invokes none of the bounded-$r$ scalar or normal conclusions
(BE5)--(BE6). Only bounded $u$ and infinitesimal $v$ remain.

\paragraph{An exact centered root identity.}
Write $D_u=1+u^2$, $D_v=1+v^2$ and introduce
\[
 h=A^{-1/2},\quad L=Ar,\qquad
 H=\sqrt A(Y-ruW),\quad B=X+uZ,\quad U=uW,\quad V=vZ.
 \tag{CAO5}
\]
Thus $Y=rU+hH$, while $h,r$ are infinitesimal and $L\sim2$.
All coefficient radicals in (CAO5) lie in the coefficient field;
$B,U,V$ lie in the testing field. The root $H$ is their finite tensor
combination, with no mixed inverse or mixed radical.

For clarity, the complete original finite parent used here is
\begin{align*}
F={}&4m(vW+uvY-X-uZ)^2+8D_v(uW-Y)^2
       +4AD_uD_vY^2\\
 &+(4D+2d)(uW+uvX-Y-vZ)^2
   +(4\alpha+2d)D_u(vY-Z)^2\\
 &+(4\beta+2d)D_v(X+uZ)^2-2dD_uD_vZ^2.
\end{align*}

Choose $0<D_*\le D$ bounded, and let $m_*=0$ or choose a positive
real $m_*\le m$, as specified below. Retain the four nonnegative
original event contributions with weights $m-m_*,D-D_*,\alpha,\beta$.
For $C=4D_*+2d$, the exact remaining parent is
\begin{align*}
F_*={}&8D_v[(1-r)U-hH]^2
 +4D_uD_vH^2+4LrD_uD_vU^2+8LhD_uD_vUH\\
&+C[(1-r)U-hH+uvB-D_uV]^2+2dD_vB^2\\
&+2dD_u\{v^2(rU+hH)^2-2(rU+hH)V-V^2\}
 +4m_*R_m^2,                                      \tag{CAO6}\\
R_m={}&vW+uv(rU+hH)-B.
\end{align*}
In particular, the centered $A$ contribution is the exact identity
\[
 4AD_uD_vY^2
 =4D_uD_vH^2+4LrD_uD_vU^2+8LhD_uD_vUH.
\]
This retains its complete hard root. The last two coefficients have
zero residue. Formula (CAO6) follows directly from the original parent
and includes the complete signed core; it is not an eliminated scalar
or a restriction to state normals. Every displayed coefficient is
bounded. The term involving $Z$ in that core has already cancelled
to the displayed expression in $V=vZ$.

\paragraph{Every state when $u$ is a bounded unit.}
Set $m_*=0$, so the entire original $m$ event is retained. Formula
(CAO6) is an exact Gram allocation on $(H,B,U,V)$. At infinitesimal
$v$ its residue is
\[
 4D_uH^2+2dB^2+8U^2+C(U-D_uV)^2-2dD_uV^2.
 \tag{CAO7}
\]
Residues are understood in this and the following determinant
calculations. The $U,V$ determinant is
\[
 D_u\{8CD_u-2d(8+C)\}.
\]
If $\bar\eta>1$, take $D_*=D$, which is bounded by (CAO4).
Then $C=8/(\bar\eta-1)$ in residue and this determinant is
\[
 \frac{64u^2(1+u^2)}{\bar\eta-1}>0.
\]
If $\bar\eta=1$, choose a sufficiently large positive real $D_*$;
this is permitted because $D$ is unbounded. The determinant is then
\[
 32(1+u^2)(D_*u^2+2u^2-2)>0.
\]
The first diagonal is $8+C>0$, so the complete four-root residue
Gram is positive definite in both cases. The cap is a real constant
chosen from the nonzero testing residue of $u$; no division of tensor
elements is introduced. The bounded full Gram is therefore tensor
SOS, and restoring the retained events proves the claim for every
state in this chart.

\paragraph{The remaining unit-$w$ state chart.}
Suppose $u,v$ are infinitesimal and $w$ is a unit. Then $W$ is a unit
on the negative-core branch. When $Z$ is infinitesimal the inherited
(BE4) argument already applies. Otherwise normalize the state so that
$Z$ is a unit and put $q=W/Z$, a pure-testing unit. Choose a positive
real lower bound $m_*\le m$. Now
\[
 R_m=qV+uv(rU+hH)-B,
\]
so (CAO6) is again an exact bounded Gram on $(H,B,U,V)$, with residue
\[
 4H^2+2dB^2+4m_*(qV-B)^2
       +8U^2+C(U-V)^2-2dV^2.                    \tag{CAO8}
\]
The $B$ pivot is positive. Its Schur complement adds
\[
 \mu V^2,\qquad \mu=\frac{4dm_*q^2}{d+2m_*}>0
\]
to the $U,V$ block. This is just a residue determinant computation;
no mixed inverse is used to define a root or a full-parent coefficient.
If $\bar\eta>1$, take $D_*=D$ as before. The determinant becomes
\[
 \frac{8\bar\eta\mu}{\bar\eta-1}>0.
\]
If $\bar\eta=1$, choose a sufficiently large positive real cap
$D_*$ so that $\mu(C+8)-64>0$. The unboundedness of $D$ permits this
choice. Thus the complete residue Gram in (CAO8) is strict, including
all hard-soft terms, and restoring the retained events gives the
full-parent tensor SOS.

\paragraph{Completion of state entry.}
Normalize the state by a maximal coordinate and first remove the
nonnegative pure-testing-core branch, which has its direct original
SOS. On the negative-core branch with $v$ infinitesimal, the residue
core is
\[
 2\{(X+uZ)^2+(uW-Y)^2\}.
\]
Consequently these two roots have zero residue. If $w$ and $Z$ were
both infinitesimal, the exact inverse relations
\[
 W=w+\frac{u(uW-Y)}{D_u},\quad
 Y=uw-\frac{uW-Y}{D_u},\quad X=(X+uZ)-uZ
\]
would make every normalized state coordinate infinitesimal. Hence
small $w$ forces $Z$ to be a unit. The bounded-unit $u$ chart was
settled by (CAO7), and all unit-$w$ states by (CAO8) and (BE4).
Normalize the only remaining states by $Z=1$. The same exact
pure-testing Schur identity (SD5) and testing sign gauges give
$u,v,w>0$. This is exactly (CAO1). Literal axes and literal state
zeros have already been included. Complex states split into their
real and imaginary evaluations because the parent is real.

The accompanying \path{check_centered_a_outer.py} independently
checks the native source cancellations, the exact selected-parent
identity, both complete root allocations, their residues and all
displayed determinants. It does not import the contact scalar
calculations used for the remaining chart.
